% Generated draft: requires mathematical and source review.
\documentclass[11pt]{article}
\usepackage[margin=25mm]{geometry}
\usepackage{fontspec}
\setmainfont{Latin Modern Roman}
\usepackage{amsmath,amssymb,mathtools,unicode-math}
\setmathfont{Latin Modern Math}
\usepackage{enumitem,needspace,xurl,hyperref}
\setlength{\parindent}{0pt}
\setlength{\parskip}{5pt}
\setlength{\emergencystretch}{3em}
\newcommand{\N}{\mathbb{N}}
\newcommand{\Z}{\mathbb{Z}}
\newcommand{\Q}{\mathbb{Q}}
\newcommand{\R}{\mathbb{R}}
\newcommand{\C}{\mathbb{C}}
\newcommand{\F}{\mathbb{F}}
\newcommand{\A}{\mathbb{A}}
\newcommand{\PP}{\mathbb{P}}
\newcommand{\E}{\mathbb{E}}
\newcommand{\eps}{\varepsilon}
\newcommand{\dd}{\,\mathrm{d}}
\begin{document}
\section*{UnsolvedMath ID: 701386}
\textbf{Source problem code:} catalogue entry
\subsection{Definitions and mathematical statement}
Let \(G\) be a group and let \(A=\C[G]\) be its complex group algebra, whose elements are finite formal sums of elements of \(G\) with complex coefficients. All tensor products below are over \(\C\). Write \(A^e=A\otimes A^{\mathrm{op}}\), so that left \(A^e\)-modules are identified with \(A\)-bimodules. Give \(A\otimes A\) the outer bimodule action
\[
a\cdot(x\otimes y)\cdot b=ax\otimes yb.
\]
For interpreting the bimodule structure on the Ext groups, use the commuting inner action
\[
a\mathbin{\star}(x\otimes y)\mathbin{\star}b=xb\otimes ay.
\]
This specifies the standard bimodule convention used in this draft; the supplied record does not state that convention explicitly.

In the terminology of the supplied source, \(A\) is \emph{Calabi--Yau of dimension \(3\)}, abbreviated CY\(3\), when
\[
\operatorname{Ext}^{3}_{A^e}(A,A\otimes A)\cong A
\quad\text{as \(A\)-bimodules},
\qquad
\operatorname{Ext}^{n}_{A^e}(A,A\otimes A)=0
\quad\text{for every integer }n\geq 0\text{ with }n\ne 3.
\]
The source calls the displayed isomorphism natural, but does not specify a naturality requirement beyond the bimodule isomorphism. Its stated definition imposes no separate homological smoothness or finite generation hypothesis; none is added here.

A group \(G\) is \emph{residually finite} if for every \(g\in G\setminus\{1\}\) there exist a finite group \(F\) and a homomorphism \(\varphi:G\to F\) such that \(\varphi(g)\ne 1\).

\textbf{Question (R. Mikhailov, 2010).} Is every group \(G\) for which \(\C[G]\) is CY\(3\) in the source's stated sense residually finite? The quantifier ranges over all groups; the supplied question imposes no finite generation or finite presentation assumption.

\textbf{Qualification concerning the supplied background.} The record attributes to Kontsevich the assertion that the complex group algebra of the fundamental group of a \(3\)-dimensional aspherical manifold is CY\(3\). It does not specify compactness, orientability, boundary conditions, or further hypotheses for that assertion, and supplies no separate citation to Kontsevich's theorem. This is retained as attributed source background, not asserted here as a theorem in that unrestricted form, and is not an additional hypothesis of the question.
\subsection{Short English statement}
R. Mikhailov asks whether every group whose complex group algebra satisfies the supplied CY\(3\) Ext conditions is residually finite: can every nonidentity element be detected by a homomorphism to a finite group? The source labels the problem open. Its Calabi--Yau terminology and its attributed manifold background require the qualifications stated above.
\subsection{Sources}
\begin{itemize}
\item Supplied local UnsolvedMath registry record, permanent ID \(701386\), exact source code \texttt{catalogue entry}, title \emph{Kourovka Notebook Problem 17.84}. Upstream revision: \texttt{d371ee917ac9647fbb06cd6fc529deec6f3dda82}. The record supplies the status \texttt{open}; this status has not been independently verified.
\item \emph{The Kourovka Notebook (Unsolved Problems in Group Theory)}, problem 17.84, proposed by R. Mikhailov in 2010; the supplied background identifies the 17th issue (2010) and the 21st issue edition (2026), \texttt{arXiv:1401.0300v48}. The supplied collection description names E. I. Khukhro and V. D. Mazurov as editors.
\item The supplied statement attributes its assertion about fundamental groups of aspherical \(3\)-manifolds to Kontsevich. No separate bibliographic reference or precise hypotheses for that theorem are supplied. The cited source has not been independently consulted for this draft.
\end{itemize}
\end{document}
